\section*{第 8 章：带类型的算术表达式}
\addcontentsline{toc}{section}{第 8 章：带类型的算术表达式}

\begin{taplsolutionentrybox}
\phantomsection
\label{sol:translator:8.2.3}
\textbf{练习 8.2.3}

证明目标是：若 \(t:T\)，则 \(t\) 的每个子项都具有某个类型。对
\(t:T\) 的类型推导作归纳。归纳命题取为：对推导中的每个结论
\(s:S\)，项 \(s\) 以及它的所有子项都是良类型的。

\textbf{基础情形。} 若最后一条规则是 \rulename{T-True}、
\rulename{T-False} 或 \rulename{T-Zero}，则 \(t\) 没有真子项；而
\(t\) 本身已由该规则获得类型，所以结论成立。

\textbf{归纳情形。} 其余每条类型规则都先给直接子项建立类型，再为整个项
建立类型。例如：

\begin{itemize}
  \item 若最后使用 \rulename{T-If}，三个前提分别说明 \(t_1,t_2,t_3\)
    良类型。由归纳假设，它们各自的所有子项也良类型；条件项的子项恰好是这
    三个直接子项及其子项，所以结论成立。
  \item 若最后使用 \rulename{T-Succ}、\rulename{T-Pred} 或
    \rulename{T-IsZero}，唯一的直接子项 \(t_1\) 由规则前提获得类型；由
    归纳假设，\(t_1\) 的所有子项也良类型。
\end{itemize}

所有类型规则均已覆盖，故每个良类型项的每个子项都是良类型的。

\taplsolutionbacklink{ex:8.2.3}{返回练习 8.2.3}
\end{taplsolutionentrybox}

\begin{taplsolutionentrybox}
\phantomsection
\label{sol:translator:8.3.4}
\textbf{练习 8.3.4}

要证：若 \(t:T\) 且 \(t\trans t'\)，则 \(t':T\)。这次对
\(t\trans t'\) 的求值推导作归纳。归纳假设是：对当前求值推导的每个直接
子推导 \(s\trans s'\)，只要 \(s:S\)，就有 \(s':S\)。按最后使用的求值
规则分类。

\begin{itemize}
  \item \rulename{E-IfTrue}：此时
    \(t=\taplif{\tapltrue}{t_2}{t_3}\) 且 \(t'=t_2\)。由 \(t:T\) 的
    反演引理，\(t_2:T\)，故结论成立。\rulename{E-IfFalse} 相同。

  \item \rulename{E-If}：此时
    \(t=\taplif{t_1}{t_2}{t_3}\)、\(t_1\trans t_1'\)，且
    \(t'=\taplif{t_1'}{t_2}{t_3}\)。由反演引理，
    \(t_1:\tapltype{Bool}\) 且 \(t_2:T\)、\(t_3:T\)。对子推导
    \(t_1\trans t_1'\) 使用归纳假设，得到
    \(t_1':\tapltype{Bool}\)；再用 \rulename{T-If} 即得 \(t':T\)。

  \item \rulename{E-Succ}、\rulename{E-Pred}、
    \rulename{E-IsZero}：都先由反演引理得到被求值子项的类型，再对子推导
    使用归纳假设，最后用相应类型规则重建 \(t':T\)。

  \item \rulename{E-PredZero}：\(t=\taplpred{\taplzero}\)、
    \(t'=\taplzero\)。由 \(t:T\) 的反演引理得到
    \(T=\tapltype{Nat}\)，而 \rulename{T-Zero} 给出
    \(t':\tapltype{Nat}\)。

  \item \rulename{E-PredSucc}：\(t=\taplpred{(\taplsucc{nv_1})}\)、
    \(t'=nv_1\)。由反演引理得到 \(T=\tapltype{Nat}\) 且
    \(nv_1:\tapltype{Nat}\)，所以 \(t':T\)。

  \item \rulename{E-IsZeroZero}：结果为 \(\tapltrue\)；
    \rulename{E-IsZeroSucc}：结果为 \(\taplfalse\)。两种情形都由反演引理
    得到 \(T=\tapltype{Bool}\)，再由常量的类型规则得到结论。
\end{itemize}

这些正是算术语言的全部求值规则，故保持性成立。

\taplsolutionbacklink{ex:8.3.4}{返回练习 8.3.4}
\end{taplsolutionentrybox}
